% Related work -- prior work grouped by scope, roughly oldest first within each group.
% Subsection of "Background", pulled in by sections/02-background.tex.

\subsection{Related work}                % prior work on airway geometry, segmentation, CFD, encoding, labelling, morphology
\label{sec:related}                      % \secref{sec:related}

% ==============================================================================
\subsubsection{Geometric models of the airway tree}  % morphometry, generation, casts, CT-based and hybrid models
\label{sec:related_geometry}             % \secref{sec:related_geometry}

\paragraph{Morphometry and idealized models.}  % from lung casts to mathematical descriptions
Early airway models are based on idealized morphometry rather than on individual anatomy.
Casts of human lungs were used to study branching patterns and the relation between airway
lengths and diameters \cite{horsfield1986}, and \citet{horsfield1971models} and
\citet{yeh1980models} proposed asymmetric models of the human bronchial tree.
\citet{kamiya1974} related airway diameter to branching angles, and \citet{wang1992} examined
deterministic bifurcating distributive systems with a Monte Carlo method.
\citet{kitaoka1999} proposed a 3D model of the human airway tree in which the branching
pattern is driven by a diameter--flow relationship within an organ boundary, relating
branching angle to flow rate and airway length to diameter. \citet{hegedus2004} gave a
detailed mathematical description of airway bifurcation geometry and used it to generate
surface models of idealized bifurcations, and \citet{florens2010} presented an anatomical and
functional model of the tracheobronchial tree. \citet{tawhai2011} reviewed computational
models of airway and pulmonary vascular structure and function towards a ``Lung Physiome''.

% \citet{varner2017computational} reviewed computational approaches to airway branching
% morphogenesis in the mammalian lung, grouping them into three classes: geometric,
% reaction--diffusion, and continuum mechanical models. Geometric models, inspired by
% Mandelbrot's fractal theory, show that simple recursive rules can generate space-filling,
% lung-like trees that match morphometric data, but they do not explain the biological
% mechanisms behind these rules. Physical models, including viscous fingering, buckling driven
% by constrained epithelial growth, and apical constriction, show that mechanical forces can
% also determine where branches form. The authors argue that each framework captures only part
% of the process, and that future models should combine biochemical and mechanical cues and be
% grounded in quantitative experimental data from embryonic lungs.

\paragraph{Algorithmic generation of airway trees.}  % volume filling and L-systems
Building on the Monte Carlo approach of \citet{wang1992}, \citet{tawhai2000} introduced a
volume-filling algorithm that grows a bifurcating tree inside each lung lobe.
\citet{davoodi2016} generated the human bronchial tree with a Lindenmayer system, and
\citet{abbasi2021} extended it to a stochastic parametric L-system with parallel growth of
bronchopulmonary segments. \citet{geronzi2024} proposed a parametric 3D model of human
airways for particle drug delivery and deposition.

% \paragraph{Lung casts and digital reference models.}  % physical replicas
% \citet{clinkenbeard2002} built hollow airway replicas of the first five generations using
% rapid prototyping. \citet{schmidt2004} created a digital reference model of the bronchial tree
% (17 generations) by segmenting a CT scan of a rubber lung cast, and \citet{lizal2012} derived
% realistic and semi-realistic airway geometries up to the seventh bifurcation from it.

\paragraph{Subject-specific models from CT.}  % CT-resolved airways, often extended algorithmically
Subject-specific airway trees can be extracted from CT images. \citet{palagyi2006} developed
the skeletonization used to obtain a 1D tree from a segmented airway, and
\citet{nakamura2008} presented automated segmentation and morphometric analysis of the airway
tree from multidetector CT. \citet{tawhai2004} combined CT-based geometry analysis with finite
element models of the human and ovine bronchial tree, using their volume-filling algorithm to
extend the CT-resolved tree to the terminal bronchioles, guided by the resolved airways and
the lobar boundaries. \citet{lin2009} and \citet{tawhai2009} combined CT-based central airways
with such artificial airways to build subject-specific models from the glottis to the
terminal bronchioles for CFD, but their pipeline required manual intervention to construct
trifurcations; \citet{lin2013} later reviewed multiscale image-based modeling of gas flow and
particle transport in the lungs. \citet{montesantos2013} reported HRCT-based airway
morphometry in healthy subjects and asthmatic patients, and \citet{montesantos2016} presented
an algorithm for an individualized 3D deterministic model of the human bronchial tree from
HRCT images, which they evaluated statistically. \citet{bordas2015} developed a pipeline for
patient-based complete conducting airway models of healthy and asthmatic subjects, covering
CT acquisition, lung, lobar and airway segmentation, centerline extraction, and algorithmic
generation of the distal airways beyond the CT-resolved tree.

Trifurcations are a frequent source of difficulty for these models, although they are common
in real airway trees. \citet{boyden1949} studied bronchopulmonary segments in 50 subjects and
reported a trifurcation at the right upper lobe bronchus in 46\% of them, and
\citet{zhao2009} examined the left-lung bronchial anatomy of 216 subjects with multi-detector
CT and found trifurcations at the left upper division bronchus and the left lower lobar
bronchus in 23\% and 18\% of subjects, respectively.

\citet{miyawaki2017automatic} proposed a centerline-based method that automatically
constructs subject-specific three-dimensional models of the human airways from a CT-segmented
airway skeleton and surface. Their approach systematically identifies and classifies
trifurcations, distinguishing true trifurcations from double bifurcations and fork-type from
tripod-type configurations, and uses dedicated geometric templates to reconstruct them. The
method produces models of increasing realism (straight-centerline, curved-centerline, and
fitted-surface), can be extended beyond CT resolution to the terminal bronchioles, and
partitions the geometry branch by branch, which facilitates automatic mesh generation for
computational fluid dynamics (CFD). Applied to 16 healthy and 16 severe asthmatic subjects,
the fitted-surface models achieved an average surface error of about 11\% of the image voxel
size, and an aerosol deposition simulation showed branch-wise deposition efficiencies within
5\% of those obtained with a CT-based model.

\citet{nousias2020avatree} presented AVATREE for generating
anatomically valid patient-specific airway tree models from CT scans. Their pipeline first
segments the lungs and the visible central airways and extracts the airway centerline as a
graph. It then extends the tree beyond the generations visible in the image using a
space-filling algorithm which splits the lung volume with PCA based planes. From the extended
centerline they reconstruct a 3D airway surface suitable for CFD simulations using Poisson
surface reconstruction. The framework also produces spatial probability maps of airway
generations overlaid on the CT images. It can simulate broncho-constriction as seen in asthma
and COPD, through Laplacian mesh contraction. The authors validated the generated trees
against published morphometric data comparing Strahler and Horsfield orders, branch lengths,
diameters and branching angles.

\paragraph{Meshing and hybrid models.}   % from segmented geometry to simulation-ready meshes
\citet{marchandise2013} proposed an automatic method to build and mesh tubular geometries,
including trifurcations, for cardiovascular and lung applications. However, it takes a surface
as input and is not centerline based, so the mesh patches are not linked to branches with
anatomical data. \citet{lauria2022} developed an automatic method to generate CFD compliant
triangulated meshes from segmented airways, \citet{ortizpuerta2022} introduced Snakes
Isogeometric Analysis based on NURBS for healthy and COPD airway trees and \citet{ho2024}
smoothed airway surface meshes with graph convolutional neural networks. Following the same
idea as \citet{bordas2015}, \citet{agujetas2020} and \citet{ilegbusi2023} built hybrid models
that couple image derived upper airways with idealized lower airways for CFD.

\citet{pennati2024modeling} presented a comprehensive review of geometrical models of the
human intrathoracic airways tracing their evolution from the early idealized mathematical
models of the 1960s such as the symmetric Weibel model and the asymmetric Horsfield model,
to modern image-based, patient-specific reconstructions derived from CT and to a lesser
extent. The authors surveyed airway segmentation strategies ranging from traditional
region-growing and knowledge-based methods to machine learning and deep learning approaches
such as 3D U-Net architectures and outlined the pipeline required to obtain CFD compliant
meshes including centerline extraction, surface smoothing and mesh generation. They also
discussed hybrid models that combine CT-resolved central airways with algorithmically
generated distal trees and reviewed applications in inter-subject variability studies,
diagnosis, surgical planning, drug delivery and environmental exposure. Finally, they
identified key open challenges including missing small airways in segmentation, limited
validation on pathological scans and the need for robust, automated mesh generation suitable
for clinical use.

% % ==============================================================================
% \subsubsection{Airway segmentation and benchmarks}  % extracting the tree from CT and MRI
% \label{sec:related_segmentation}         % \secref{sec:related_segmentation}

% \paragraph{Traditional methods.}         % region growing, rules, morphology, tubes
% Early methods segmented the airway tree from CT by region growing \cite{mori1995}, later
% combined with rule-based and knowledge-based strategies exploiting the relation between
% airways and vessels \cite{sonka1996}, fuzzy logic \cite{park1998}, mathematical morphology
% and topology \cite{pisupati1996,fan2000}, tube detection \cite{bauer2009tube} and gradient
% vector flow \cite{bauer2009gvf}. \citet{lo2009} extracted airway trees via locally optimal
% paths, and \citet{irving2011} segmented obstructed airway branches using airway topology and
% statistical shape analysis. \citet{smistad2014} introduced a GPU-accelerated segmentation and
% centerline extraction framework for tubular structures, and later the FAST framework for
% heterogeneous medical image computing and visualization \cite{smistad2015}. \citet{pu2012}
% reviewed CT-based identification and analysis of human airways.

% \paragraph{Machine learning and deep learning.}  % voxel classification to CNNs
% Voxel-classification approaches with kNN classifiers \cite{lo2008,lo2010} and AdaBoost with
% minimum spanning trees \cite{inoue2013} reduced leakage and improved the detection of small
% branches. Deep learning approaches include CNN-based leak detection \cite{charbonnier2017},
% 2.5D CNNs \cite{yun2019} and 3D U-Net architectures \cite{garciauceda2021}, as well as
% freeze-and-grow propagation combined with deep learning \cite{nadeem2021}, tubule-sensitive
% CNNs \cite{qin2021}, bronchiole-sensitive networks with attention distillation
% \cite{qin2020}, context scale fusion \cite{qin2020airwaynet}, and neural networks with linear
% programming \cite{zhao2019}. \citet{tan2022} reviewed the segmentation of lung airways with
% deep learning.

% \paragraph{Benchmarks and datasets.}     % public data for training and comparison
% The EXACT'09 challenge \cite{lo2009exact09,lo2012exact} compared fifteen airway extraction
% algorithms on a common set of CT scans. More recent datasets provide larger benchmarks: ATM'22
% \cite{zhang2023atm22}, AIIB23 \cite{nan2024aiib23} and AeroPath \cite{stoverud2023aeropath}
% cover multi-site data, fibrotic lung disease and challenging pathology, respectively. ATM'22
% and AIIB23 are the datasets used in this work (\secref{sec:datasets}).

% \paragraph{MRI-based reconstruction.}    % airways without ionising radiation
% \citet{lewis2005} and \citet{peterson2011} used hyperpolarized $^3$He MRI to measure airway
% diameters and reconstruct the airway tree in three dimensions, \citet{wang2004} proposed
% scale-based fuzzy connectedness for 3D airway segmentation on hyperpolarized gas MRI, and
% \citet{genkin2024} segmented airways from the trachea to the tertiary bronchi with ultrashort
% echo time MRI in pediatric patients.

% ==============================================================================
% \subsubsection{Airway geometries for flow simulation}  % CFD and particle deposition
% \label{sec:related_cfd}                  % \secref{sec:related_cfd}

% \paragraph{Idealized and early computational models.}  % geometry built from morphometric models
% The idealized models of \citet{horsfield1971models} and \citet{yeh1980models} were used by
% \citet{vanertbruggen2005} and \citet{tian2011sip} to build centerline-based 3D airway
% geometries for CFD and particle deposition. These models are not subject-specific and do not
% include trifurcations. \citet{gemci2008} and \citet{walters2011} also built computational
% airway models for airflow simulation, but without an automatic construction procedure, and
% \citet{jiang2020} proposed a voxel image-based airflow simulation with automatic detection of
% the airway outlets.

% \paragraph{Patient-specific CFD and particle deposition.}  % CT-based geometries in simulation
% CT-based patient-specific airway geometries have been used for CFD studies of airflow,
% including steady inspiration in reconstructed proximal airway trees \cite{vial2005}, flow
% analysis in the lower airways with patient-specific boundary conditions \cite{debacker2008},
% bifurcating flow in a CT-scanned airway \cite{luo2008}, airflow with subject-specific boundary
% conditions \cite{yin2010}, and airflow and particle transport in generations three to six
% \cite{srivastav2011}. \citet{swan2012} modelled the topographic distribution of ventilation in
% healthy lungs. Other works address aerosol distribution in a CT-based airway model
% \cite{miyawaki2012}, pharmaceutical aerosol deposition validated against in vivo data
% \cite{tian2015}, static versus dynamic imaging for particle transport
% \cite{miyawaki2016static}, large-scale CFD of airflow and particle deposition \cite{soni2013},
% and direct numerical simulation of particle-laden flow in an airway bifurcation
% \cite{stylianou2016}. \citet{kolanjiyil2016} developed a computationally efficient
% whole-lung-airway model for particle transport and deposition. Studies of obstructive
% conditions and further flow effects include normal versus obstructed airways \cite{sul2014},
% dynamic flow in normal and asthmatic lungs \cite{kim2015}, the turbulent laryngeal jet
% \cite{lin2007}, 4DCT-based breathing lung models \cite{miyawaki2016dynamic}, and airflow in
% the acinar region \cite{kumar2009}. \citet{nousias2016} and \citet{lalas2017} simulated
% broncho-constriction through geometry contraction of CT-extracted airways and used it to
% assess substance deposition in obstructed lungs.

% ==============================================================================
\subsubsection{Representation learning and generative models of vascular trees}  % encoding / synthesis
\label{sec:related_representation}       % \secref{sec:related_representation}

Data-driven approaches have recently emerged as an alternative to rule-based (fractal and
space-filling) methods for synthesizing 3D vascular geometry. \citet{feldman2023vesselvae}
proposed VesselVAE, a recursive variational neural network (RvNN) that encodes vessels as
binary trees where each node stores a centerline point and its radius. The encoder
recursively aggregates the tree into a low-dimensional latent space while the decoder
reconstructs node attributes and predicts through a node classifier, whether each node
branches. New vessels are obtained by sampling the latent space and converting the generated
centerlines into surface meshes. Trained on the IntrA dataset the method produced samples
whose radius, length and tortuosity distributions closely match real data. However, it
assumes circular cross-sections and has difficulty scaling to deep trees. To overcome these
limitations, the same group introduced VesselGPT \cite{feldman2025vesselgpt}, which treats
vascular trees as sequences of nodes obtained through a preorder traversal. A VQ-VAE first
learns a discrete vocabulary of node embeddings and a GPT-2 model is then trained to generate
vessels autoregressively over these tokens. Instead of a single radius each cross-section is
represented with a B-spline which preserves non-circular shapes such as aneurysms. The final
mesh is extracted from a signed distance field using marching cubes. Evaluated on the
Aneurisk dataset, VesselGPT generates deeper and more detailed vascular trees than previous
recursive and diffusion-based approaches. Graph-based generation has also been explored:
\citet{prabhakar2024vesselgraph,prabhakar2025semantic} proposed two-stage discrete diffusion
models that generate simplified metric graphs of vessels and airways.

Learning compact representations of vascular structures has recently gained attention with
approaches differing mainly in how the vessel geometry is represented.
\citet{batten2025vetta} proposed VeTTA, a two-stage Transformer-based autoencoder that encodes
vessel trees into a single vector representation. In the first stage a Vessel Autoencoder
embeds individual vessel segments described by centerline points and radii lifted to
sinusoidal Fourier features into compact latent codes. In the second stage a Vessel Tree
Autoencoder encodes the full set of tree edges into a global latent vector which is decoded
recursively with a slot-based Transformer decoder that guarantees a topologically valid tree.
Evaluated on a synthetic 2D tree dataset and a 3D coronary artery dataset, VeTTA outperformed
3D convolutional autoencoders in reconstruction accuracy required less GPU memory and
produced smooth, topologically consistent latent interpolations. However, it is restricted to
tree-structured data with at most two children per node.

In contrast, \citet{prabhakar2026vaesselsparse} proposed VAEsselSparse a variational
autoencoder that operates directly on high-resolution organ-level binary vessel segmentation
maps without assuming a tree structure. By combining sparse 3D convolutions with sparse
windowed self-attention their model exploits the inherent sparsity of vascular masks and
achieves an $8\times8\times8$ spatial compression while preserving vessel connectivity.
VAEsselSparse outperformed a dense VAE and MedVAE on topology-aware metrics such as clDice and
Betti number errors. The authors further showed that the learned latent space supports
downstream tasks including aneurysm/stenosis classification, circle of Willis subtype
classification and the generation of realistic vessel structures with a flow-matching model
trained in the latent space.

% ==============================================================================
\subsubsection{Branch matching and anatomical labeling of airway trees}  % correspondence and labels
\label{sec:related_labeling}             % \secref{sec:related_labeling}

Registration, branch matching and branch labeling of airway trees have been studied for over
a decade. Several successful methods are based on association graphs where branch matchings
are induced by maximal cliques of a graph combining information from two trees
\cite{kitaoka2002automated,tschirren2005matching,metzen2009matching};
\citet{tschirren2005matching} reported a labeling success rate of 97\%.
\citet{vanginneken2008robust} labeled airway branches recursively starting at the trachea
as part of the segmentation process assigning labels from radius, orientation and
bifurcation angle (90\% success rate); this approach is vulnerable to differences in tree
topology. \citet{kaftan2006novel} matched tree paths rather than branches which avoids
topological differences but discards the information stored in the tree structure and does
not yield branch labels. \citet{bulow2006point} matched airway branches using only branch
shape without connectivity information (69\% or 40\% success rate depending on the features
used).

In contrast to these methods which focus on either topology or branch geometry
\citet{feragen2011airway} proposed a geodesic tree-shape model for anatomical labeling of
airway branches that treats tree topology and branch geometry jointly. Each airway centerline
tree is represented as a combinatorial tree whose edges carry landmark points sampled along
the branches. Trees of different sizes and topologies are embedded in a common Euclidean
space and different representations of the same shape are identified giving a quotient
space equipped with the Quotient Euclidean Distance (QED). In this space any two trees are
connected by a shortest deformation (geodesic) that can change the tree topology along the
way. Since computing exact geodesics is NP-hard, the authors approximate them by limiting the
number of topological changes in each lobar subtree. Branch correspondences obtained from the
geodesics between an unlabeled tree and a set of expert-labeled trees are combined through
majority voting to assign anatomical labels. 
% Evaluated in a leave-one-out setting on the 20
% scans of the EXACT'09 test set \cite{lo2009exact09}, segmented with the voxel-classification
% method of \citet{lo2010}, the method correctly labeled 83\% of the branches on average,
% excluding one outlier case with missing upper lobes. 
\citet{feragen2012hierarchical} later
turned these geodesic distances into a hierarchical labeling scheme: a new airway tree is
labeled from the root toward the periphery by matching it against a set of manually labeled
training trees.

\citet{weng2026topofield} proposed TopoField, a topology-aware implicit modeling framework for
repairing incomplete pulmonary trees (airways, arteries, and veins) extracted from CT scans.
The pulmonary tree is represented by sparse surface and skeleton point clouds which are fused
through a surface-to-skeleton cross-attention module and projected onto a triplane implicit
field that can be queried at arbitrary spatial locations. To avoid the need for annotated
breakpoints the authors introduce TopoBreak a strategy that synthetically adds
disconnections to already-incomplete trees and trains the network to recover them providing
weak supervision for topology repair. On top of the shared implicit field, task-specific
decoders jointly perform tree repair, anatomical labeling and lung segment reconstruction in
a single forward pass. 
% This approach achieves accurate results on the Lung3D+ benchmark while
% running in about one second per case, roughly two orders of magnitude faster than previous
% implicit graph-based pipelines.

% ==============================================================================
\subsubsection{Morphological and topological analysis of airway trees}  % quantifying and comparing shape
\label{sec:related_morphology}           % \secref{sec:related_morphology}

\paragraph{Local airway geometry.}       % wall thickness and lumen shape
% Automatic extraction of the airway tree from CT is a prerequisite for morphological analysis
% (\secref{sec:related_segmentation}).
Early quantitative work focused on local airway
geometry. \citet{nakano2000airway} related the wall area percent of a segmental bronchus and
the extent of emphysema to lung function in smokers, and \citet{hasegawa2006airflow} extended
such measurements across airway generations using multiplanar reconstruction.
\citet{patel2008pi10} used Pi10, a standardized measure of airway wall thickness, to show that
wall thickening and emphysema aggregate independently in families. \citet{oguma2015irregularity}
quantified the irregularity of the lumen radius along a single airway path and found it
higher in COPD than in healthy subjects and asthma.

\paragraph{Global metrics of the tree.}  % one number per tree
Several global metrics have been derived from the segmented tree. \citet{kirby2018tac}
introduced the total airway count (TAC) and showed that fewer visible branches predict COPD
progression. \citet{bodduluri2018afd} characterized the complexity of the tree by its fractal
dimension and linked lower values to morbidity and mortality in COPD, and later proposed the
airway surface-area-to-volume ratio as a marker of airway remodeling associated with airflow
limitation \cite{bodduluri2021sav}. \citet{smith2020dysanapsis} defined the airway-to-lung
size ratio (dysanapsis) and showed that a smaller ratio is associated with COPD risk, while
\citet{maetani2023dysanapsis} studied the physiological effects of this ratio, fractal
dimension, and branch count in healthy never smokers.

\paragraph{Statistics on tree shapes.}   % comparing whole trees, including topology
\citet{feragen2011means} modeled anatomical trees as points in a space of tree-like shapes
that captures both branch geometry and topological structure, and studied how to compute mean
trees in this space, with airway trees as the application. \citet{sorensen2011dissimilarity}
represented airway trees as sets of attributed branches, matched the branches across
subjects and classified COPD using a dissimilarity-based embedding with a $k$-NN classifier.
\citet{feragen2013geometric} proposed a family of geometric tree kernels for comparing
anatomical trees whose branches carry vector-valued attributes such as branch shape, radius
or airway wall area percentage. Unlike classical graph and tree kernels which rely on
discrete node labels their path-based (all-pairs and root-path) and pointcloud kernels
incorporate continuous geometric information while remaining efficient enough to scale to
datasets of nearly 10{,}000 trees. Applying these kernels to airway trees segmented from
low-dose CT scans of 1966 subjects they used kernel-based two-sample tests to show that
airway trees of COPD patients and healthy individuals follow significantly different
distributions. They also showed using SVM classification that geometry-aware kernels
outperform average airway wall area percentage measures and state of the art graph kernels
such as the shortest path and Weisfeiler - Lehman kernels in detecting COPD.

\paragraph{Topological descriptors.}     % persistent homology
\citet{belchi2018topology} applied persistent homology to the airway tree using height as the
filtration and showed that the resulting ``upwards complexity'' correlates with COPD
severity. \citet{kaji2024airway} proposed two morphological metrics based on persistent homology to
quantify tubular structures in chest CT images and applied them to the airway trees of
patients with chronic obstructive pulmonary disease (COPD). They computed 0-dimensional
persistent homology under two distance based filtrations. The first, the geodesic distance
from the carina yields \emph{treeH$_0$}, which describes the complexity of the branching
structure. The second the distance from the lumen centerline yields \emph{radialH$_0$},
which captures irregularities of the luminal surface in the central (1st to 5th generation)
airways. On three cohorts acquired with different scanners and reconstruction kernels, both
metrics were higher in COPD patients than in non-COPD smokers. In addition, radialH$_0$
predicted longitudinal FEV$_1$ decline independently of conventional CT measures such as
emphysema extent (LAA\%), airway-to-lung size ratio, and total airway count.
